\textbf{Phần 1. Trắc nghiệm lựa chọn}
\setcounterpageref{ex}{0}

\begin{ex} %[2H6N29]
Một ngày có bao nhiêu giờ?
\choice
{$12$ giờ}
{$20$ giờ}
{\True $24$ giờ}
{$60$ giờ}
\loigiai{
Một ngày có 24 giờ.\\
\textbf{Đáp án của câu hỏi là C}
}
\end{ex}

\begin{ex} %[2H6N29]
Một giờ có bao nhiêu phút?
\choice
{\True $60$ phút}
{$30$ phút}
{$100$ phút}
{$24$ phút}
\loigiai{
Một giờ có 60 phút.\\
\textbf{Đáp án của câu hỏi là A}
}
\end{ex}

\begin{ex} %[2H6H29]
$1$ ngày được tính từ lúc
\choice
{$6$ giờ sáng đến $6$ giờ tối}
{$7$ giờ sáng đến $12$ giờ trưa}
{$1$ giờ sáng đến $12$ giờ đêm}
{\True $12$ giờ đêm hôm trước đến $12$ giờ đêm hôm sau (từ 0 giờ đến 24 giờ)}
\loigiai{
1 ngày tính từ 12 giờ đêm hôm trước đến 12 giờ đêm hôm sau.\\
\textbf{Đáp án của câu hỏi là D}
}
\end{ex}

\begin{ex} %[2H6H29]
Lúc 8 giờ tối còn gọi là
\choice
{$18$ giờ}
{$21$ giờ}
{\True $20$ giờ}
{$22$ giờ}
\loigiai{
$8 + 12 = 20$ giờ.\\
\textbf{Đáp án của câu hỏi là C}
}
\end{ex}

\begin{ex} %[2H6H29]
Lúc 15 giờ còn gọi là
\choice
{$2$ giờ chiều}
{\True $3$ giờ chiều}
{$4$ giờ chiều}
{$5$ giờ chiều}
\loigiai{
$15 - 12 = 3$ giờ chiều.\\
\textbf{Đáp án của câu hỏi là B}
}
\end{ex}

\begin{ex} %[2H6H29]
Kim ngắn chỉ số 8, kim dài chỉ số 12. Lúc đó là
\choice
{$12$ giờ}
{$8$ giờ $30$ phút}
{\True $8$ giờ đúng}
{$7$ giờ $55$ phút}
\loigiai{
Đồng hồ chỉ 8 giờ đúng.\\
\textbf{Đáp án của câu hỏi là C}
}
\end{ex}

\begin{ex} %[2H6V29]
Kim ngắn chỉ giữa số 3 và số 4, kim dài chỉ số 6. Lúc đó là
\choice
{$3$ giờ $15$ phút}
{\True $3$ giờ $30$ phút (hoặc 3 giờ rưỡi)}
{$4$ giờ $30$ phút}
{$3$ giờ $6$ phút}
\loigiai{
Kim dài chỉ số 6 là 30 phút. Lúc đó là 3 giờ 30 phút.\\
\textbf{Đáp án của câu hỏi là B}
}
\end{ex}

\begin{ex} %[2H6V29]
Lan bắt đầu học bài lúc 7 giờ tối và học xong lúc 8 giờ 30 phút tối. Lan đã học trong thời gian là
\choice
{\True $1$ giờ $30$ phút}
{$1$ giờ}
{$2$ giờ}
{$30$ phút}
\loigiai{
Thời gian học: 8 giờ 30 phút - 7 giờ = 1 giờ 30 phút.\\
\textbf{Đáp án của câu hỏi là A}
}
\end{ex}

\begin{ex} %[2H6V29]
Buổi sáng em vào lớp lúc 7 giờ 30 phút, tan học lúc 10 giờ 30 phút. Em học ở trường trong số giờ là
\choice
{$2$ giờ}
{$4$ giờ}
{$5$ giờ}
{\True $3$ giờ}
\loigiai{
$10 - 7 = 3$ giờ.\\
\textbf{Đáp án của câu hỏi là D}
}
\end{ex}

\begin{ex} %[2H6V29]
Bố đi công tác trong 3 ngày. Bố đi công tác trong số giờ là
\choice
{$48$ giờ}
{\True $72$ giờ}
{$60$ giờ}
{$96$ giờ}
\loigiai{
$24 \times 3 = 72$ giờ.\\
\textbf{Đáp án của câu hỏi là B}
}
\end{ex}

\begin{ex} %[2H6V29]
Một tiết học kéo dài 40 phút. Nếu bắt đầu lúc 8 giờ thì tiết học kết thúc lúc
\choice
{$8$ giờ $30$ phút}
{$8$ giờ $50$ phút}
{$9$ giờ}
{\True $8$ giờ $40$ phút}
\loigiai{
Kết thúc lúc 8 giờ 40 phút.\\
\textbf{Đáp án của câu hỏi là D}
}
\end{ex}

\begin{ex} %[2H6C29]
Bây giờ là 9 giờ sáng. Hỏi sau 15 giờ nữa sẽ là mấy giờ đêm?
\choice
{\True $12$ giờ đêm (24 giờ)}
{$11$ giờ đêm}
{$1$ giờ đêm}
{$10$ giờ đêm}
\loigiai{
$9 + 15 = 24$ giờ (tức 12 giờ đêm).\\
\textbf{Đáp án của câu hỏi là A}
}
\end{ex}

\textbf{Phần 2. Trắc nghiệm đúng sai}
\setcounterpageref{ex}{0}

\begin{ex} %[2H6H29]
Xét tính đúng sai về thời gian:
\choiceTF
{\True Một ngày có 24 giờ.}
{\True Một giờ có 60 phút.}
{14 giờ là 4 giờ chiều.}
{\True 14 giờ là 2 giờ chiều ($14 - 12 = 2$).}
\loigiai{
\begin{itemize}
    \item 1 ngày = 24 giờ: Đúng.
    \item 1 giờ = 60 phút: Đúng.
    \item 14 giờ = 2 giờ chiều (không phải 4 giờ). Ý 3 Sai, Ý 4 Đúng.
\end{itemize}
\begin{itemize}
    \item a) \textbf{Đúng}.
    \item b) \textbf{Đúng}.
    \item c) \textbf{Sai}.
    \item d) \textbf{Đúng}.
\end{itemize}
}
\end{ex}

\begin{ex} %[2H6V29]
Mai xem phim từ 19 giờ đến 20 giờ 30 phút. Xét tính đúng sai:
\choiceTF
{\True 19 giờ chính là 7 giờ tối.}
{\True Bộ phim kéo dài 1 giờ 30 phút.}
{20 giờ 30 phút chính là 9 giờ 30 phút tối.}
{\True 20 giờ 30 phút chính là 8 giờ 30 phút tối ($20 - 12 = 8$).}
\loigiai{
\begin{itemize}
    \item 19 giờ = 7 giờ tối: Đúng.
    \item Thời gian: 1 giờ 30 phút: Đúng.
    \item 20 giờ = 8 giờ tối (không phải 9 giờ). Ý 3 Sai, Ý 4 Đúng.
\end{itemize}
\begin{itemize}
    \item a) \textbf{Đúng}.
    \item b) \textbf{Đúng}.
    \item c) \textbf{Sai}.
    \item d) \textbf{Đúng}.
\end{itemize}
}
\end{ex}

\textbf{Phần 3. Trả lời ngắn}
\setcounterpageref{ex}{0}

\begin{ex} %[2H6H29]
Một ngày có bao nhiêu giờ?
\shortans{24}
\loigiai{
\textbf{Đáp số:} $24$\\
24 giờ.
}
\end{ex}

\begin{ex} %[2H6H29]
Một giờ có bao nhiêu phút?
\shortans{60}
\loigiai{
\textbf{Đáp số:} $60$\\
60 phút.
}
\end{ex}

\begin{ex} %[2H6V29]
Lúc 17 giờ chính là mấy giờ chiều?
\shortans{5}
\loigiai{
\textbf{Đáp số:} $5$\\
$17 - 12 = 5$ giờ chiều.
}
\end{ex}

\begin{ex} %[2H6V29]
Lúc 9 giờ tối tương ứng với bao nhiêu giờ trong ngày?
\shortans{21}
\loigiai{
\textbf{Đáp số:} $21$\\
$9 + 12 = 21$ giờ.
}
\end{ex}

\begin{ex} %[2H6V29]
Nam tập thể dục từ 6 giờ sáng đến 6 giờ 30 phút sáng. Nam đã tập thể dục bao nhiêu phút?
\shortans{30}
\loigiai{
\textbf{Đáp số:} $30$\\
30 phút.
}
\end{ex}

\begin{ex} %[2H6C29]
2 ngày có tất cả bao nhiêu giờ?
\shortans{48}
\loigiai{
\textbf{Đáp số:} $48$\\
$24 + 24 = 48$ giờ.
}
\end{ex}
